%%%%%%%%%%%%%%%%%%%%%%%%%%%%%%%%%%%%%%%%%%%%%%%%%%%%%%%%%%%%%%%%%%%%
%% appendix1.tex
%% ISEL thesis document file
%%
%% Chapter with example 
%%%%%%%%%%%%%%%%%%%%%%%%%%%%%%%%%%%%%%%%%%%%%%%%%%%%%%%%%%%%%%%%%%%%
\chapter{Apêndice 1} 
\label{apendice:1}

\section{Casos de utilização}
\label{apendice:casos_util}

Nesta secção serão apresentados os casos de utilização do projeto. Todos os casos de utilização têm apenas um ator, sendo este o utilizador da aplicação.

\subsection{Inserção de dados através de um ficheiro}
Resumo: Leitura de um ficheiro, que siga o padrão do concurso \glsxtrshort{itc} 2019, e armazenamento dos dados numa base de dados.

Cenário principal:

\begin{enumerate}
    \item O caso de utilização inicia-se quando o utilizador pressiona o botão de inserção de dados através de ficheiro.
    \item O utilizador escolhe o ficheiro que deve ser lido por parte da aplicação. \label{casos_util_1:escolha_ficheiro}
    \item A aplicação processa os dados e armazena-os na base de dados. \label{casos_util_1:processar_dados}
    \item O caso de utilização termina.
\end{enumerate}

Cenário alternativo 1:

\begin{enumerate}
    \item No passo~\ref{casos_util_1:escolha_ficheiro} o utilizador cancela a escolha do ficheiro.
    \item O caso de utilização termina.
\end{enumerate}

Cenário alternativo 2:

\begin{enumerate}
    \item No passo~\ref{casos_util_1:processar_dados} o ficheiro escolhido não segue o padrão do concurso \glsxtrshort{itc} 2019.
    \item Apresentação de um erro ao utilizador.
    \item O caso de utilização termina.
\end{enumerate}

\subsection{Inserção de dados manual}
Resumo: Inserção de dados relativos a salas, professores ou disciplinas manualmente.

Cenário principal:

\begin{enumerate}
    \item O caso de utilização inicia-se quando o utilizador pressiona o botão de inserção de dados manual (este pode ser relativo a salas, professores ou disciplinas).
    \item O utilizador preenche a informação relativa à categoria escolhida. \label{casos_util_2:preench_dados}
    \item A aplicação processa os dados recebidos e armazena-os na base de dados. \label{casos_util_2:processar_dados}
    \item O caso de utilização termina.
\end{enumerate}

Cenário alternativo 1:

\begin{enumerate}
    \item No passo~\ref{casos_util_2:preench_dados} o utilizador pressiona o botão para cancelar a inserção de dados.
    \item O caso de utilização termina.
\end{enumerate}

Cenário alternativo 2:

\begin{enumerate}
    \item No passo~\ref{casos_util_2:processar_dados} os dados introduzidos não estão de acordo com o esperado.
    \item Apresentação de uma mensagem de erro ao utilizador.
    \item O caso de utilização termina.
\end{enumerate}

\subsection{Criação de horários}
Resumo: Após o clique no botão de criação de horários, o programa deve gerar os horários, conforme a configuração especificada, e apresentar os horários gerados na interface gráfica.

Pré-condições: Devem haver dados armazenados na base de dados.

Cenário principal:

\begin{enumerate}
    \item O caso de utilização inicia-se quando o utilizador pressiona o botão de criação de horários.
    \item A aplicação processa os dados armazenados e cria os horários pedidos.
    \item São apresentados os horários gerados através da interface gráfica.
    \item O caso de utilização termina.
\end{enumerate}

\subsection{Visualização de um horário já criado}
Resumo: Após selecionado um dos horários gerados, de uma lista de horários, apresentados na interface gráfica, deve ser apresentado o horário de forma ampliada.

Pré-condições: Devem ter sido gerados horários previamente.

Cenário principal:

\begin{enumerate}
    \item O caso de utilização inicia-se quando o utilizador seleciona um horário gerado de uma lista de horários gerados pela aplicação.
    \item A aplicação apresenta o horário de forma ampliada ao utilizador.
    \item O caso de utilização termina.
\end{enumerate}

\subsection{Modificação de um horário criado}
Resumo: Após selecionado um dos horários gerados, de uma lista de horários, apresentados na interface gráfica, deve ser apresentado o horário permitindo a modificação do mesmo.

Pré-condições: Devem ter sido gerados horários previamente.

Cenário principal:

\begin{enumerate}
    \item O caso de utilização inicia-se quando o utilizador seleciona um horário gerado de uma lista de horários gerados pela aplicação.
    \item A aplicação apresenta o horário de forma ampliada ao utilizador.
    \item O utilizador modifica o horário ao mover as atribuições das aulas.
    \item O utilizador confirma as modificações através do clique num botão. \label{casos_util_5:confirmar_alteracoes}
    \item A aplicação armazena as alterações num novo horário de modo a não se perder o horário original.
    \item O caso de utilização termina.
\end{enumerate}

Cenário alternativo 1:

\begin{enumerate}
    \item No passo~\ref{casos_util_5:confirmar_alteracoes} o utilizador pressiona o botão para cancelar as alterações efetuadas ao horário.
    \item O caso de utilização termina.
\end{enumerate}

\subsection{Exportação dos horários criados}
Resumo: São exportados os horários selecionados no formato escolhido para a localização que o utilizador especificou.

Pré-condições: Devem ter sido gerados horários previamente.

Cenário principal:

\begin{enumerate}
    \item O caso de utilização inicia-se quando o utilizador seleciona a opção de exportar horários.
    \item O utilizador seleciona o formato no qual o horário deve ser exportado, sendo os seguintes formatos suportados \glsxtrshort{xml}, \glsxtrshort{pdf}, \glsxtrshort{csv} e \glsxtrshort{png}.
    \item O utilizador deve apresentar o local para o qual os horários serão exportados.
    \item A aplicação efetua a exportação do horário.
    \item O caso de utilização termina.
\end{enumerate}